% Fragmento público generado desde propietarios íntegros.
% Residencia: 52.11.1.
% Fuente: ../incoming/radion_monografia_20260827/manuscrito/sections/12_aritmetica_borde.tex:1-159
\noindent The nonadic boundary preserves the sum--product defect, vacancy, and completion carry.\par\medskip

\subsubsection{Radical \(3\! -\! 6\! -\! 9\) of the boundary}

\paragraph{Generation and nilpotency of the nonadic radical}

In the ring \(\Z/9\Z\), the ideal
\begin{equation}
\mathfrak m=(3)=\{0,3,6\}
\label{eq:radical-369}
\end{equation}
satisfies
\begin{equation}
\mathfrak m^2=(9)=0.
\label{eq:m-square-zero}
\end{equation}
It is the nilpotent sector of the nonadic boundary. The visible orbit of its three
marks is published as
\[
369\longrightarrow693\longrightarrow936\longrightarrow369.
\]
The return does not imply that all states are equal: the visible
permutation closes while the route cocycles may advance.

The double convention
\[
9\equiv0\pmod9,
\qquad
\operatorname{dr}_9(9)=9
\]
expresses two readers: algebraic residual zero and visible digital-root
closure. Thus a nine may absorb boundary torsion without becoming
a specific zero of the zeta function.

\paragraph{Sum--product defect \(459-405=54\)}

The APP sheets produce complementary totals
\[
S_+=405,\qquad S_\times=459,\qquad S_\times-S_+=54.
\]
The \(54\) is an exact defect of the tables and reappears as a contraction
exponent through the identity \(9\cdot6=54\). Each occurrence nevertheless
requires its own map. The focal distance \(d_1=0.544545\ldots\) is not equal to
\(0.54\), and invariant analysis has produced no arrow identifying it with
the \(54\) defect.

The power
\[
9^9=387,420,489
\]
generates a scale where
\begin{equation}
\operatorname{ord}_{10}\bigl(9^{9^9}\bigr)=369,693,099,
\qquad
\#\,\text{digits}=369,693,100.
\label{eq:369693}
\end{equation}
The block \(369|693\) appears in the order--length pair; it is not asserted to be
the decimal prefix of the power.

\subsubsection{Arithmetic Theory of Decimal Vacuums}

\paragraph{Asymptotic vacancy slope}

The separation between the decade and the nonary base is
\begin{equation}
\delta_{\mathrm{vac}}=\log_{10}\frac{10}{9}.
\label{eq:delta-vac}
\end{equation}
The accumulated counter
\begin{equation}
V_9(n)=\left\lceil n\delta_{\mathrm{vac}}\right\rceil
\label{eq:vacancy-count}
\end{equation}
records how many correction positions are required by the decimal publication of
a nonadic evolution. Its jumps form a binary word
\[
z_n=V_9(n)-V_9(n-1)\in\{0,1\}.
\]
The gaps are distributed with lengths \(21/22\), the returns with
lengths \(6/7\), and the combined determinant is \(15\). These numbers
arise from the vacancy reader, not from the Barbero kernel, although \(15\)
coincides with \(\binom62\).

\paragraph{Polarization and Arithmetic Current of the Nonadic Grid}

We define
\begin{equation}
P_N=\sum_{n=1}^N(z_n-\delta_{\mathrm{vac}}),
\qquad
J_N=P_N-P_{N-1}=z_N-\delta_{\mathrm{vac}}.
\label{eq:arith-current}
\end{equation}
The balanced sequence satisfies \(|P_N|<1\). The \enquote{current}
\(J_N\) is here an arithmetic current: it measures the creation and absorption of
vacancies relative to the mean slope. Only a subsequent dimensional transducer
can publish it as an electromagnetic current density.

The physical intuition is powerful because the same algebraic form separates an
even average from an oriented fluctuation. Rigor requires retaining the type:
\[
J_N\in\R\quad\text{dimensionless},
\qquad
j^\mu\in\text{physical current line}.
\]

\subsubsection{Completion \(729\to1000\) and preservation of carry}

\paragraph{Decomposition \(1000=729+243+27+1\)}

The completion
\[
1000=729+243+27+1
\]
preserves the unit through four ternary scales. The total crown is
\[
1000-729=271,
\]
while the last integer before the closure satisfies
\[
999=729+270.
\]
The difference between \(270\) and \(271\) distinguishes an occupied boundary from a
completing unit. The same care prevents identifying a vacancy with zero or a
carry with an actual residue.

\paragraph{Residence of the carry of the radion in the decimal crown}

The radion's operands begin, in blocks of three digits, with
\[
\Phi_T:\quad000|027|455|906|837|565|\cdots,
\]
\[
D_E:\quad000|027|455|010|034|743|\cdots.
\]
The subtraction APP gives
\[
\epsone:\quad000|000|000|896|802|822|\cdots.
\]
The small size of the result is not obtained by truncating the low-order
summands. It arises because two distinct publications share a long prefix,
and the operator \(\operatorname{SubAcarreo}_{1000}\) preserves the borrows required to subtract
them exactly.

The word \emph{memory} has two levels here:
\begin{enumerate}
\item \(q_{\mathrm{mem}}\) records the depth of nonadic returns;
\item \(c_i\) records the local borrowings of decimal subtraction.
\end{enumerate}
The first organizes the helix; the second publishes the decimal. \(\epsR\) uses
both because its operands follow from cylinders at depth and its subtraction is
performed with carry, but it is equal to neither counter.

\begin{thesisbox}
The boundary is not a place where the theory loses information. It is the
place where information changes representation: phase that returns, memory
that ascends, cylinder that contracts, vacancy that is recorded, and carry
that transports the difference.
\end{thesisbox}
